\subsection{Controlled follow-up with Gemma 12B}\label{sec:followup}
\paragraph{CPU prefix-size ablation.}
We varied only the persona prefix supplied to the Agent constructor in the existing long-persist fixture: the current prefix, committed-HEAD prefix, removed prefix, and neutral text matched to the current prefix in character count and UTF-8 byte lengths. Three interleaved repetitions per arm shared one fixed temporary path. The provider remained scripted; its tokenizer was synthetic ByteLevel BPE with no merges, context 32,768 and output reserve 512. All arms generated 37 requests with identical non-system character counts at every round.

Final message-character totals were 11,344 (current), 10,598 (HEAD), 9,886 (removed), and 11,344 (neutral), identical across the three repetitions. Only the removed-prefix arm passed the unchanged threshold, in 3/3 repetitions. The current final request contained 4,693 system, 1,077 user, 47 assistant and 5,527 tool characters. The 1,458-character prefix exactly explained its difference from the removed-prefix arm; a neutral prefix reproduced the size result without preserving persona semantics. The HEAD arm also failed, so the latest uncommitted persona edit is not the sole explanation. The earlier 11,364 observation used a different temporary path and is retained as a distinct measurement.

This identifies a size-sensitive assertion and an accounting boundary: the history reducer does not impose a universal character cap on system/evidence plus all protected/latest content. It does not measure persona quality or justify removing persona instructions in deployment. Maximum observed current-arm input was 32,014 synthetic tokens; adding its eight-token safety allowance and output reserve remained below 32,768. Thus this fixture's character failure did not imply token-admission failure, and these synthetic counts are not Gemma counts.

\paragraph{Gemma deployment and compatibility.}
The owner selected Gemma 12B for this stage. We tested local alias \texttt{gemma4:12b}, Q4\_K\_M, Ollama 0.33.3, manifest \texttt{4eb23ef187e2} on the same laptop. Tokenizer/config/template assets from Google's instruction-tuned repository \cite{gemmaassets} were pinned to revision \texttt{43c9d014b231}. Local GGUF metadata names the base repository; correspondence to that upstream instruction-tuned weight revision is not independently established. The local manifest and asset hashes define the tested deployment. No production model setting was changed.

Eight bounded provider requests covered Ukrainian text, emoji/code-switching, a tool schema and completed tool history, twice each. Context was 4096, output reserve 32, temperature zero, seed 42 and thinking disabled. Every returned provider input count exactly matched the local count: 32, 27, 75 and 98 by configuration. All ended with stop; proposed toy tools were never executed. This is a narrow calibration, excluding near-window boundaries and long sessions.

\paragraph{Current-code task pilot.}
We then ran three hand-selected existing cases: Ukrainian factual answering, file write/read, and memory storage followed by recall in a new session. Two repetitions at contexts 4096 and 8192 yielded 12 executions and 16 user turns. All 12 passed the existing deterministic scorer: 6/6 at each context. Output reserve was 128, temperature zero, top\_p 1, seed 42, with maximum eight model rounds, twelve tool calls and a 240-second experiment budget. The existing runner/scorer and runtime sources were unchanged; an experiment wrapper supplied explicit Gemma settings because the CLI defaults to Bolt. Each execution used fresh disposable stores and sandbox. Both bounded model probes acknowledged unloading their alias afterward.

These observations establish selected current-code execution paths with an additional model, not a reliability estimate or superiority over Bolt. Three selected scenarios and correlated repetitions cannot support ranking; the historical Bolt run also used different code, corpus coverage and output reserve. Neither shared-model-switch state nor blinded behavioral identity was tested. Full corpus, working-note and provenance ablations remain future experiments. Scripts, registered protocol, source/corpus hashes, raw checks, JUnit and cleanup records accompany this stage.

